
\section{Discussão}

Os resultados obtidos demonstram que a utilização de técnicas de agrupamento possibilita organizar estudantes em grupos com características acadêmicas semelhantes, permitindo uma análise mais estruturada dos padrões de desempenho presentes na base de dados. Essa capacidade de identificação de perfis representa uma etapa importante para o desenvolvimento de sistemas inteligentes de apoio à decisão acadêmica, uma vez que a caracterização dos estudantes pode futuramente subsidiar estratégias de recomendação personalizadas.

A avaliação da incorporação de Grandes Modelos de Linguagem permitiu definir uma arquitetura compatível com as restrições computacionais e financeiras da pesquisa. Embora tais modelos apresentem potencial para enriquecer a interação com os usuários, a análise demonstrou que sua adoção não era adequada às condições atuais do projeto, motivando a priorização de uma solução baseada em algoritmos tradicionais de Mineração de Dados.

\subsection{Limitações}

Apesar dos resultados obtidos demonstrarem a viabilidade da aplicação de técnicas de Mineração de Dados para a caracterização de perfis acadêmicos, algumas limitações devem ser consideradas na interpretação deste estudo.

A principal limitação está relacionada à base de dados inicialmente prevista para a pesquisa. Os experimentos seriam realizados a partir dos históricos acadêmicos dos estudantes de Ciência da Computação da Universidade Federal Rural do Rio de Janeiro (UFRRJ). Entretanto, durante a etapa de processamento dos documentos, foram identificadas inconsistências decorrentes da formatação dos históricos gerados pelo sistema acadêmico, comprometendo a extração estruturada e confiável das informações necessárias para os experimentos. Em razão disso, foi necessário utilizar bases de dados acadêmicas alternativas para validar a metodologia de tratamento dos dados e a aplicação do algoritmo de agrupamento. Consequentemente, os resultados apresentados não podem ser generalizados para os estudantes da UFRRJ, mas representam uma validação inicial da abordagem proposta em dados estruturados.

Outra limitação refere-se ao estágio de desenvolvimento da arquitetura. Embora o objetivo final do projeto seja oferecer recomendações personalizadas de disciplinas e auxiliar o planejamento da trajetória acadêmica dos estudantes, o módulo de recomendação ainda se encontra em fase de implementação. Assim, esta etapa da pesquisa concentra-se na identificação e caracterização dos perfis acadêmicos por meio do algoritmo K-Means, não sendo possível avaliar, neste momento, o impacto das recomendações sobre o desempenho ou sobre a redução do tempo de formação.


